%!TEX root = ../csuexperiment_main.tex

\newcommand{\reportchapter}[1]{%
    \chapter*{#1}%
    \phantomsection%
    \addcontentsline{toc}{chapter}{#1}%
}
\newcommand{\reportsection}[1]{\section*{#1}}

\reportchapter{实验目的}

本实验旨在围绕实验主题建立完整的问题理解、方法实现与结果分析流程。通过完成实验，掌握相关理论方法的基本原理，熟悉实验环境配置、程序实现、结果记录和误差分析等步骤，并能够结合输出结果对实验现象进行解释。

\reportchapter{实验原理}

本章说明实验涉及的核心概念、关键算法、系统结构或数学模型。写作时应从基本定义出发，逐步解释方法成立的原因、主要公式的含义以及实验方案与理论之间的对应关系。

\reportchapter{实验内容与步骤}

\reportsection{（一）实验环境}

本实验的运行环境、开发工具、依赖库和主要参数如下。生成正式报告时，应结合实际运行环境补充操作系统、编程语言版本、编译器或解释器版本以及关键第三方库版本。

\reportsection{（二）实验步骤}

本实验按照问题分析、环境配置、程序实现、结果记录和结果分析的顺序展开。正式报告中应写清每一步的输入、处理过程、输出文件和关键判断依据。

\reportchapter{结果与分析}

\reportsection{（一）运行结果}

本节展示实验运行得到的主要结果，包括关键输出、统计表、对比图或程序运行截图。每个图表都应在正文中先说明来源，再结合结果解释其含义。

\reportsection{（二）结果分析}

本节从正确性、性能、稳定性、边界情况和误差来源等角度分析实验结果。若结果与预期存在差异，应说明可能原因并给出验证或改进思路。

\reportchapter{心得体会}

通过本次实验，对相关理论、实现过程和结果分析方法有了更完整的理解。正式报告中应结合实验中的具体问题，说明调试过程、关键收获、仍可改进的方面以及后续可以继续探索的方向。

\reportchapter{参考文献}

\begin{enumerate}
    \item 待补充与实验主题相关的教材、论文、官方文档或课程资料。
\end{enumerate}
